\section*{Chapitre \ref{ch:s1:running-a-multi-strategy-book} --- Gérer un portefeuille multi-stratégies}

\begin{solution}{exo:s1:running-a-multi-strategy-book:1}
La position nette de dix transactions indépendantes et égales a un écart-type égal à $\sqrt{10}$ fois celui d'une seule, contre un brut de dix~: le netting économise $1 - 1/\sqrt{10} = 68.4~\%$
(la simulation mesure 68.5~\%).
\end{solution}

\begin{solution}{exo:s1:running-a-multi-strategy-book:2}
$0.4 \times 6 = 2.4$ volatilités, soit 24~\% du capital à une volatilité de 10~\%.
\end{solution}

\begin{solution}{exo:s1:running-a-multi-strategy-book:3}
Leur corrélation se mesure surtout sur des journées calmes, où le facteur commun n'est qu'un faible bruit~; dans un krach, le facteur commun domine et les trois évoluent
avec lui. Une corrélation moyenne faible masque une importante exposition commune dans la queue de distribution.
\end{solution}

\begin{solution}{exo:s1:running-a-multi-strategy-book:4}
34.5~\%~: à 5~\% de volatilité, une limite de 5~\% vaut une volatilité annuelle, comme une limite de 10~\% à 10~\% de volatilité. Ce qui compte, c'est la limite mesurée dans les
volatilités propres du \omterm{def:s1:running-a-multi-strategy-book:multi}{pod}.
\end{solution}

\begin{solution}{exo:s1:running-a-multi-strategy-book:5}
Tout \omterm{def:s1:running-a-multi-strategy-book:multi}{pod}, mort ou vivant, atteint rapidement une limite serrée~; la dérive négative du \omterm{def:s1:running-a-multi-strategy-book:multi}{pod} mort l'y conduit plus tôt (98 jours à 5~\%). Le coût est d'arrêter presque tous les \omterm{def:s1:running-a-multi-strategy-book:multi}{pods} vivants
chaque année (91.6~\% des trajectoires), en perdant leur edge jusqu'à leur redémarrage.
\end{solution}

\begin{solution}{exo:s1:running-a-multi-strategy-book:6}
Si la richesse reste au-dessus d'une fraction de son maximum courant, la position risquée optimale est proportionnelle au surplus au-dessus de ce plancher~: le risque
diminue progressivement à mesure que les pertes approchent du plancher et croît à mesure que le surplus se reconstitue. Diviser par deux le risque au premier drawdown et l'arrêter au second est une approximation en deux
étapes de cette droite.
\end{solution}

\begin{solution}{exo:s1:running-a-multi-strategy-book:7}
Avec un cap égal au quart d'un \omterm{def:s1:running-a-multi-strategy-book:multi}{pod}, les pertes lors des krachs tombent à 12.2~\%, 11.1~\% et 13.0~\%, le plus grand drawdown à 28.0~\% et le ratio de Sharpe monte à 1.72.
Dans une vraie firme, la carte des facteurs est estimée, non connue~: un cap strict sur des facteurs mesurés avec erreur réduirait à tort des \omterm{def:s1:running-a-multi-strategy-book:multi}{pods} pour des expositions qu'ils n'ont pas, et
pousserait les \omterm{def:s1:running-a-multi-strategy-book:multi}{pods} vers des expositions que le modèle de risque ne voit pas.
\end{solution}

\begin{solution}{exo:s1:running-a-multi-strategy-book:8}
Trois années de données, surtout calmes, mesurent la corrélation du bruit, non celle des expositions qui comptent en stress~; les \omterm{def:s1:running-a-multi-strategy-book:multi}{pods} qui se recouvraient dans la firme synthétique affichaient
0.22 de corrélation en moyenne et perdirent ensemble un quart de leur capital. Il faut mesurer les expositions à partir des positions et tester la firme dans des scénarios de stress.
\end{solution}

\begin{solution}{pb:s1:running-a-multi-strategy-book:1}
\begin{enumerate}
\item Un portefeuille combiné de stratégies gérées par des équipes~; une équipe disposant de son propre capital et de ses propres limites~; la corrélation de leurs P\&L, qui peut masquer
des expositions communes.
\item Dix \omterm{def:s1:running-a-multi-strategy-book:multi}{pods} à 10~\% de volatilité, ratios de Sharpe de 0.5 à 1.0~; \omterm{def:s1:running-a-multi-strategy-book:multi}{pods} 1--3 et 4--5 sur deux facteurs communs avec trois krachs~; \omterm{def:s1:running-a-multi-strategy-book:multi}{pod} 10 mort après dix ans.
\item Par inverse de la volatilité glissante, mensuellement.
\item Un désendettement généralisé du marché dans les stratégies quantitatives sur actions, avec des débouclages les 1 et 6 août.
\item 0.22 entre \omterm{def:s1:running-a-multi-strategy-book:multi}{pods} sur un même facteur, $-0.004$ avec les autres~; 2.4 à 2.6 volatilités perdues ensemble à chaque krach.
\item À partir des positions projetées sur les facteurs du modèle de risque, avec un plafond sur l'exposition totale de chaque facteur.
\item 1.45, 43.5~\% et 31.5--33.6~\% sans~; 1.70, 35.7~\% et 20.7--22.9~\% avec un cap de la moitié d'un \omterm{def:s1:running-a-multi-strategy-book:multi}{pod}.
\item 34.6~\% lorsque les \omterm{def:s1:running-a-multi-strategy-book:multi}{pods} négocient un cinquième des instruments, 68.5~\% lorsque tous négocient la totalité des instruments.
\item Un arrêt sur une perte depuis le sommet~; arrêter les \omterm{def:s1:running-a-multi-strategy-book:multi}{pods} vivants et laisser tourner les \omterm{def:s1:running-a-multi-strategy-book:multi}{pods} morts.
\item 91.6~\%, 34.5~\%, 8.9~\%, 2.1~\% et 0.0~\% sur les trajectoires~; 85.8~\%, 38.9~\%, 22.1~\%, 14.7~\% et 5.8~\% dans la firme, pour 5~\% à 30~\%.
\item 98, 307 et 979 jours à 5~\%, 10~\% et 15~\%~; jamais à 20~\% et 30~\%.
\item Un seul échantillon, cinq candidats, et des écarts inférieurs au bruit.
\item \textbf{Résultat nommé.} À 10~\% de volatilité de la firme, le plus grand drawdown atteint 43.5~\% sans détection des recouvrements et 35.7~\% lorsque les facteurs communs sont plafonnés
à la moitié d'un \omterm{def:s1:running-a-multi-strategy-book:multi}{pod}~; une \omterm{def:s1:running-a-multi-strategy-book:limit}{limite de drawdown} de 5~\% arrête 91.6~\% des \omterm{def:s1:running-a-multi-strategy-book:multi}{pods} sains (ratio de Sharpe 0.7, volatilité de 10~\%) en moins d'un an.
\item La position risquée optimale est proportionnelle au surplus au-dessus d'un plancher fixé comme une fraction du maximum courant.
\item Dans les volatilités propres au \omterm{def:s1:running-a-multi-strategy-book:multi}{pod} et le ratio de Sharpe attendu.
\item Plusieurs \omterm{def:s1:running-a-multi-strategy-book:multi}{pods} perdent ensemble, chacun à cause de sa propre limite.
\item Positions et leurs expositions aux facteurs, P\&L quotidiens, et transactions pour le netting.
\item Le livre~7 mesurait l'encombrement (crowding) entre firmes~; ce chapitre le mesure au sein d'une seule.
\item Les \omterm{def:s1:running-a-multi-strategy-book:multi}{pods} les plus exposés au facteur qui a chuté, ramenés au cap~; les autres n'ont pas causé la perte.
\item Une vue du risque commun entre \omterm{def:s1:running-a-multi-strategy-book:multi}{pods}, du netting et des limites fixées dans les unités propres à chaque \omterm{def:s1:running-a-multi-strategy-book:multi}{pod}.
\end{enumerate}
\end{solution}

\begin{solution}{iq:s1:running-a-multi-strategy-book:1}
À partir de la volatilité attendue et du ratio de Sharpe du \omterm{def:s1:running-a-multi-strategy-book:multi}{pod}~: fixer la limite de sorte que la probabilité d'arrêter un \omterm{def:s1:running-a-multi-strategy-book:multi}{pod} doté de l'edge attendu dans l'année soit
acceptable (quelques pour cent), avec une étape souple qui réduit d'abord le risque de moitié~; la réviser à mesure que l'historique du \omterm{def:s1:running-a-multi-strategy-book:multi}{pod} s'étoffe.
\end{solution}

\begin{solution}{iq:s1:running-a-multi-strategy-book:2}
Une exposition commune~: projeter les positions des deux portefeuilles sur des facteurs communs (styles encombrés, liquidité, une même contrepartie ou une même source de financement), vérifier si les deux
détiennent les mêmes valeurs ou transactions sous des libellés différents, et examiner le positionnement connu d'autres firmes.
\end{solution}

\begin{solution}{iq:s1:running-a-multi-strategy-book:3}
Comparer la performance récente à la distribution implicite du ratio de Sharpe attendu, chercher un mécanisme qui aurait pu la détruire (encombrement, coûts,
régime), et admettre qu'une baisse du ratio de Sharpe de 0.9 à $-0.3$ demande environ un an ou plus pour être détectée avec confiance.
\end{solution}

\begin{solution}{iq:s1:running-a-multi-strategy-book:4}
Collecter les ordres des \omterm{def:s1:running-a-multi-strategy-book:multi}{pods} par instrument sur une courte fenêtre, effectuer un croisement interne des ordres opposés à un prix de référence équitable, n'envoyer au marché que le solde net, et
répartir équitablement les exécutions et les coûts entre les \omterm{def:s1:running-a-multi-strategy-book:multi}{pods}~; conserver une piste d'audit.
\end{solution}

\begin{solution}{iq:s1:running-a-multi-strategy-book:5}
Réduire le risque avant que la limite ne le fasse, en gardant les positions à meilleur edge~; la firme doit souhaiter une réduction progressive du risque à mesure que le drawdown croît,
plutôt qu'un \omterm{def:s1:running-a-multi-strategy-book:multi}{pod} qui joue pour se refaire ou un \omterm{def:s1:running-a-multi-strategy-book:multi}{pod} qui liquide tout à la limite.
\end{solution}

\begin{solution}{iq:s1:running-a-multi-strategy-book:6}
Pour $\mu = 0$, $P(D_T > L) = P(|W_T| > L) = 2\big(1 - \Phi(L/(\sigma\sqrt{T}))\big)$~; pour $L = 0.5\sigma$ et $T = 1$ an, elle vaut 0.617. Le drawdown peut
dépasser $L$ à un instant antérieur à $T$ et se résorber d'ici $T$, si bien que la probabilité qu'il dépasse $L$ à un moment de l'intervalle avant $T$ est plus grande~; c'est ce sur quoi
se déclenche une règle d'arrêt.
\end{solution}
